\section{DiT, ConvNeXt y la antifragilidad}

Si el capítulo anterior trataba sobre todo de metodología, ConvNeXt y DiT muestran cómo esa metodología se concreta en obras específicas. Lo más interesante de ConvNeXt no es que «la red convolucional haya vuelto a ganar», sino que se atrevió a plantear una pregunta más incisiva justo cuando el Vision Transformer estaba en su apogeo: ¿por qué funciona realmente ViT? Al colaborar con el pasante Liu Zhuang (刘壮), Xie Saining no se limitó a hacer remiendos nostálgicos sobre la red convolucional; en cambio, trasladó de vuelta al marco convolucional toda una serie de hábitos de entrenamiento y diseño que la era de ViT había demostrado eficaces, y luego analizó, elemento por elemento, cuáles factores importaban de verdad. La conclusión final es sumamente provocadora: la soft attention quizá no sea la parte más determinante de ViT; muchos de los beneficios provienen de una receta de entrenamiento más general, de la normalización, del estilo de patchify y del macro design.

\begin{knowledgebox}{El verdadero valor de ConvNeXt}
ConvNeXt no se limita a «demostrar que la convolución no ha muerto»; a través de un conjunto de ablations sumamente convincente, desmonta los múltiples componentes del mito de ViT. Recuerda a los investigadores que, cuando un paradigma tiene un gran éxito, uno de los trabajos más importantes no es imitarlo ciegamente, sino identificar qué componentes son los que realmente actúan y cuáles son solo accesorios propios del estilo de la época.
\end{knowledgebox}

Esta postura de atreverse a cuestionar la explicación dominante se manifiesta de forma aún más clara en DiT. El punto de partida de DiT no fue «quiero publicar un artículo de gran éxito sobre modelos generativos», sino una pregunta bastante básica: ¿cómo se compara la representación que aprende un diffusion model con la del aprendizaje autosupervisado? En el proceso de investigación descubrieron que la primera era muy inferior a la segunda, y entonces pensaron: puesto que la U-Net ya se ha convertido en la opción predeterminada dentro del pipeline de diffusion, ¿por qué no usar directamente ViT como backbone denoiser? El resultado fue muy simple y muy eficaz. ¿Qué tan simple? Tan simple que, al enviar el artículo a CVPR, fue rechazado, y una de las razones fue «es demasiado simple». Este episodio parece casi una ironía de la historia de la investigación sobre sí misma: muchas ideas verdaderamente potentes son rechazadas precisamente por ser demasiado directas, por tener demasiado poco adorno, y no encajar con la imagen que los revisores tienen de la «innovación compleja».

\begin{importantbox}{DiT: la fuerza de una estructura simple}
El núcleo de DiT no está en cuántos módulos nuevos introduce, sino en que capta el juicio más determinante de una era de escalamiento: si el Transformer ya ha demostrado su unificación y su escalabilidad en otras modalidades, entonces el marco de diffusion también debería admitir un backbone más simple y más escalable. Cuanto más cerca se está de la estructura fundamental, más posible es empujar el sistema hacia una escala mayor y una interpretación más clara.
\end{importantbox}

Más tarde, DiT se envió a otra conferencia y fue aceptado como oral, y LeCun incluso publicó un tuit específicamente para burlarse de los revisores. A partir de esta experiencia, Xie Saining pronuncia otra frase que vale la pena recordar: \textit{``que un research paper sea aceptado o no, en absoluto importa... es un proceso completamente aleatorio y puro''}. Esto no significa, por supuesto, que la revisión carezca de todo valor, sino que, en las zonas límite de la investigación de vanguardia, la distinción entre muchos trabajos ya es de por sí extremadamente sutil, y el resultado de la revisión está altamente influido por el gusto particular del revisor, la presión de tiempo y su conocimiento previo. Si se eleva el hecho de ser aceptado o no a un juicio sobre el propio valor, es muy fácil entregar el ritmo de la investigación a un sistema que, por naturaleza, tiene mucho ruido. Más concretamente: cuando nació DiT, FAIR ya había comenzado su transformación cultural, y después de que él se fue, ni siquiera se le permitió figurar como autor, lo cual demuestra, además, que un trabajo verdaderamente influyente no siempre obtiene una «atribución perfecta» en el momento.

\begin{warningbox}{Tomar el resultado de la revisión como la verdad es una autolesión común entre los investigadores}
Tanto DiT como SiT pasaron por el camino de «ser rechazados primero y reconocidos después». Lo que Xie Saining quiere subrayar no es que él haya sufrido una injusticia, sino que el investigador debe aceptar un hecho: el sistema de revisión solo puede muestrear de forma tosca y no puede medir por completo el valor de un trabajo. Si se lo confunde con un veredicto final, se acabará desperdiciando una gran cantidad de tiempo en explicaciones emocionales, en lugar de seguir avanzando.
\end{warningbox}

Más adelante, DiT entró en el contexto de OpenAI a través de Bill Peebles y terminó teniendo un parentesco evidente con Sora, lo que lo hace parecer más un artículo situado en la frontera entre la investigación y la industria. Xie Saining incluso lo cuenta como solo «0.25 artículos» que verdaderamente cambiaron el rumbo de la IA, y esto no significa menospreciarse a sí mismo, sino reconocer con extrema lucidez que muchos trabajos, cuando una gran tendencia ya se ha formado, simplemente empujan un poco más allá ese límite. Precisamente por tener este sentido de la escala es que después pudo proponer la palabra clave «antifragilidad». SiT sigue un camino similar: la combinación de flow matching y Transformer volvió a ser rechazada primero y aceptada después, y estos impactos aleatorios no destruyeron la trayectoria de la investigación, sino que, al contrario, ampliaron la influencia que el trabajo tuvo después.

\begin{importantbox}{La investigación debe ser un sistema antifrágil}
La definición que ofrece Xie Saining es muy práctica: si el beneficio que trae un random shock es mayor que la pérdida, ese sistema es antifrágil. Para la investigación, esto significa que no se puede depositar todo el valor en un único envío, en una sola afiliación institucional o en una sola línea de trabajo. Al contrario, un buen diseño de investigación debe permitir que el fracaso, el retraso, el error de juicio e incluso el rechazo tengan la oportunidad de convertirse en una difusión, una reorganización o un redescubrimiento aún mayores.
\end{importantbox}

La historia de la demo temprana de Perplexity también suele citarse como hindsight comedy: cuando Aravind mostró la demo en una cafetería de Blue Bottle, Xie Saining pensó «esto no es más que GPT con un cascarón encima» y declinó la invitación. Lo verdaderamente interesante de este episodio no es demostrar que alguien se equivocó al juzgar, sino mostrar que el futuro suele formarse a partir de la superposición de muchos juicios de baja dimensión. Hay quienes dejan pasar oportunidades evidentes de conversión en producto, y también quienes juzgan mal el valor a largo plazo de una investigación; lo importante no es acertar siempre, sino que la propia estructura de capacidades no dependa de acertar en cada ocasión. El pensamiento antifrágil vuelve a ser válido aquí: hay que construir una vida y un sistema de investigación tales que, incluso si de vez en cuando se juzga mal, no queden definidos por un único error.

\subsection{Resumen de la sección}

ConvNeXt y DiT muestran juntos que lo que Xie Saining hace mejor no es seguir lo que está de moda, sino, dentro de un paradigma dominante, seguir preguntando «cuál es la variable que realmente está actuando». Junto con una conciencia lúcida de la aleatoriedad de la revisión y el énfasis en los sistemas antifrágiles, este capítulo muestra una postura de investigación capaz de soportar impactos y convertirlos en ganancia.
